%% =====================================================================
%% sec3_parameters.tex -- Model parameter accounting (Section III insert)
%% Generated 2026-07-29. Answers the advisor's question: "which parameters
%% are actually estimated, and how does each enter the dynamics?"
%%
%% USAGE (two options):
%%   (a) Drop-in replacement for the "Dynamics and parameterization" +
%%       "MHE" paragraphs of \section{System Overview and Problem Setup}
%%       in main.tex (currently lines ~148-158):  %% =====================================================================
%% sec3_parameters.tex -- Model parameter accounting (Section III insert)
%% Generated 2026-07-29. Answers the advisor's question: "which parameters
%% are actually estimated, and how does each enter the dynamics?"
%%
%% USAGE (two options):
%%   (a) Drop-in replacement for the "Dynamics and parameterization" +
%%       "MHE" paragraphs of \section{System Overview and Problem Setup}
%%       in main.tex (currently lines ~148-158):  %% =====================================================================
%% sec3_parameters.tex -- Model parameter accounting (Section III insert)
%% Generated 2026-07-29. Answers the advisor's question: "which parameters
%% are actually estimated, and how does each enter the dynamics?"
%%
%% USAGE (two options):
%%   (a) Drop-in replacement for the "Dynamics and parameterization" +
%%       "MHE" paragraphs of \section{System Overview and Problem Setup}
%%       in main.tex (currently lines ~148-158):  %% =====================================================================
%% sec3_parameters.tex -- Model parameter accounting (Section III insert)
%% Generated 2026-07-29. Answers the advisor's question: "which parameters
%% are actually estimated, and how does each enter the dynamics?"
%%
%% USAGE (two options):
%%   (a) Drop-in replacement for the "Dynamics and parameterization" +
%%       "MHE" paragraphs of \section{System Overview and Problem Setup}
%%       in main.tex (currently lines ~148-158):  \input{sec3_parameters}
%%   (b) Standalone preview:  latexmk -pdf preview_sec3.tex
%%
%% REQUIRES in the preamble (all already present in main.tex):
%%   amsmath, booktabs, siunitx, multirow, array
%%
%% NOTE ON NOTATION: main.tex currently writes the inertia increment as
%% "dJ" (\S III, \S IV-E, \S VI-C). This file uses \dJ -> \Delta J, which
%% is unambiguous against the differential d. If you adopt it, sweep the
%% remaining occurrences: grep -n "dJ" main.tex
%% =====================================================================

% ---- Notation macros (safe to move to the preamble) -----------------
\providecommand{\dJ}{\Delta J}
\providecommand{\cvec}{\mathbf{c}}
\providecommand{\rp}{\mathbf{r}_p}
\providecommand{\mpay}{\hat{m}_p}          % payload-mass scalar (prior)
\providecommand{\mb}{m_b}                % bare-airframe mass
\providecommand{\mt}{m_t}                % composite mass
\providecommand{\Tphys}{T_{\mathrm{phys}}}
\providecommand{\tauphys}{\boldsymbol{\tau}_{\mathrm{phys}}}
\providecommand{\tcom}{\boldsymbol{\tau}_{\mathrm{com}}}

% =====================================================================
% TABLE I -- parameter identity ledger
% Place early in Section III. Uses table* (two-column span); if space is
% tight, switch to table and drop the "Value / source" column.
% =====================================================================
\begin{table*}[t]
\caption{Model parameters: which are estimated online, which are algebraically
derived, and which are fixed. \emph{Only the mass is a decision variable of the
estimator}; the inertia increment and CoM offset enter the MHE as \emph{known}
runtime parameters, on the same footing as the measured thrust and torque.}
\label{tab:params}
\centering
\footnotesize
\begin{tabular}{@{}l p{0.17\textwidth} c p{0.25\textwidth} p{0.28\textwidth}@{}}
\toprule
Symbol & Meaning & Dim. & Identity & Enters the model through \\
\midrule
$m$ & composite mass (airframe + payload) & 1 & \textbf{MHE decision variable} ($\dot m = 0$, state 14) & translational dynamics \eqref{eq:trans}; cost hover baseline \eqref{eq:hover} \\
$\dJ$ & roll/pitch inertia increment & 1 & algebraically derived, \eqref{eq:dj} & rotational dynamics \eqref{eq:rot}; inner-loop gain scaling \eqref{eq:gain} \\
$\cvec = [c_x, c_y]$ & composite-CoM horizontal offset & 2 & derived \eqref{eq:cxy}, or algebraically inverted online \eqref{eq:conline} & thrust eccentricity torque \eqref{eq:tcom}; cost torque baseline \eqref{eq:hover} \\
\midrule
$\mpay$ & payload-mass scalar (\emph{the single unknown behind $\dJ$ and $\cvec$}) & 1 & weak operational prior (nominal parcel mass) & \eqref{eq:dj}, \eqref{eq:cxy} \\
$\rp = [r_x, r_y, r_z]$ & payload offset in body frame & 3 & known (gripper contact geometry) & \eqref{eq:dj}, \eqref{eq:cxy} \\
\midrule
$\mb$ & bare-airframe mass & 1 & fixed calibration & \SI{2.0643}{\kg} \\
$J_b$ & bare-airframe inertia & 3 & fixed calibration & $\mathrm{diag}(0.0142, 0.0142, 0.0210)$ \si{\kg\meter\squared} \\
$g$ & gravity & 1 & fixed constant & \SI{9.81}{\meter\per\second\squared} \\
\bottomrule
\end{tabular}
\end{table*}

% =====================================================================
% Replacement for the "Dynamics and parameterization" paragraph
% =====================================================================
\textbf{Dynamics and parameterization.} With state
$x = [\mathbf{p}; \mathbf{v}; \mathbf{q}; \boldsymbol{\omega}] \in \mathbb{R}^{13}$ and input
$u = [T; \boldsymbol{\tau}] \in \mathbb{R}^{4}$ expressed in an ENU world frame, and with
$\mathbf{e}_3 = [0, 0, 1]^{\!\top}$ the vertical basis vector, the model is
\begin{align}
\dot{\mathbf{p}} &= \mathbf{v}, \label{eq:kin}\\
\dot{\mathbf{v}} &= \frac{1}{m}\, R(\mathbf{q})\,\mathbf{e}_3 T - g\,\mathbf{e}_3, \label{eq:trans}\\
\dot{\boldsymbol{\omega}} &= J(\dJ)^{-1}\Bigl( \boldsymbol{\tau} + \tcom(\cvec) - \boldsymbol{\omega} \times J(\dJ)\boldsymbol{\omega} \Bigr), \label{eq:rot}
\end{align}
with the attitude kinematics, for $\mathbf{q} = [q_w, q_x, q_y, q_z]^{\!\top}$,
\begin{equation}
\dot{\mathbf{q}} = \frac{1}{2}
\begin{bmatrix}
-q_x & -q_y & -q_z \\
 q_w & -q_z &  q_y \\
 q_z &  q_w & -q_x \\
-q_y &  q_x &  q_w
\end{bmatrix} \boldsymbol{\omega},
\label{eq:quat}
\end{equation}
and $J(\dJ) = \mathrm{diag}\bigl(J_{xx} + \dJ,\; J_{yy} + \dJ,\; J_{zz}\bigr)$. Thrust acts at
the geometric center along the body $z$ axis, so about a composite CoM at
$\mathbf{r}_c = [c_x, c_y, c_z]$ it produces
\begin{equation}
\tcom(\cvec) = (-\mathbf{r}_c) \times \mathbf{e}_3 T = \bigl[-c_y T,\; +c_x T,\; 0\bigr]^{\!\top}. \label{eq:tcom}
\end{equation}
Note that $c_z$ \emph{cannot} appear: a vertical CoM offset is collinear with the body-$z$
thrust and their cross product vanishes identically, and gravity acts at the CoM itself. The
CoM offset is therefore a two-parameter quantity by construction, not by truncation.

Each parameter has exactly one point of entry: $m$ appears only as the thrust-to-acceleration
gain $1/m$ in \eqref{eq:trans}, $\dJ$ only in $J(\dJ)$ in \eqref{eq:rot}, and $\cvec$ only in
\eqref{eq:tcom}. This is what makes $m$ observable from a model-independent thrust
reconstruction: given $\Tphys$, the measured $\dot{\mathbf{v}}$ pins $m$ uniquely.

\textbf{Structured parameterization of the payload geometry.} $\dJ$ and $\cvec$ are
\emph{not} additional estimated degrees of freedom. Modelling the payload as a point mass
$\mpay$ rigidly attached at a known body-frame offset $\rp$, and writing
$\mt = \mb + \mpay$ and the reduced mass $\mu = \mb\,\mpay / \mt$, the parallel-axis theorem and
the CoM definition give them in closed form:
\begin{align}
\dJ &= \tfrac12\bigl(\dJ_{xx} + \dJ_{yy}\bigr) = \mu\Bigl( r_z^2 + \tfrac12\bigl(r_x^2 + r_y^2\bigr) \Bigr), \label{eq:dj}\\
\cvec &= \frac{\mpay}{\mt}\,\bigl[r_x,\; r_y\bigr]. \label{eq:cxy}
\end{align}
Three numbers are thus generated from a \emph{single} unknown scalar $\mpay$ plus known
geometry: the rigid-body constraint replaces two degrees of freedom that a naive
augmentation would have to identify. For a \SI{0.3}{\kg} payload at
$\rp = [0, 0.109, -0.47]\,$\si{\meter} this yields $\mu = \SI{0.262}{\kg}$,
$\dJ = \SI{0.0594}{\kg\meter\squared}$ --- an inertia ratio $(J_{xx}+\dJ)/J_{xx} = 5.2$ ---
and $c_y = \SI{13.8}{\milli\meter}$, i.e.\ a constant eccentric hover torque
$|c_y| \mt g \approx \SI{0.32}{\newton\meter}$, about \SI{64}{\percent} of the available
torque authority. Compensating mass alone is therefore insufficient.

\textbf{MHE.} The estimator augments the mass into the state,
$x_{\mathrm{aug}} = [x;\, m] \in \mathbb{R}^{14}$ with $\dot m = 0$ (constant over the window)
and process noise on the 13 physical states only; the payload geometry is passed in as a
\emph{known} runtime parameter vector alongside the measured input,
\begin{equation}
p_{\mathrm{MHE}} = \bigl[\, T,\; \boldsymbol{\tau},\; \dJ,\; c_x,\; c_y \,\bigr] \in \mathbb{R}^{7}. \label{eq:pmhe}
\end{equation}
Window $N = 20$, $dt = \SI{0.1}{s}$, terminal measurement excluded from the cost, and $m$
bounded to $[\SI{1.961}{\kg}, \SI{5.0}{\kg}]$ (a physical floor at $0.95\mb$: the vehicle
cannot be lighter than the bare airframe).

\textbf{NMPC parameter consumption.} The controller carries
$p_{\mathrm{NMPC}} = [x_{\mathrm{ref}}(13),\, m,\, \dJ,\, \cvec(2)] \in \mathbb{R}^{17}$.
The hover baseline in the input cost is re-centred by the live parameters,
\begin{equation}
T_h = m\,g, \qquad
u_{\mathrm{hover}} = \bigl[\, T_h,\; c_y T_h,\; -c_x T_h,\; 0 \,\bigr]^{\!\top}, \label{eq:hover}
\end{equation}
the torque entries being exactly $-\tcom$ of \eqref{eq:tcom}: hovering with an eccentric load
\emph{requires} a constant trim torque, and a cost that does not admit this pulls the torque
toward zero, producing the same class of steady-state bias as a stale hover-thrust baseline.
Finally, model consistency cannot restore inner-loop bandwidth, so the PX4 rate gains are
scaled at attach by the inertia ratio
\begin{equation}
K \leftarrow K_{\mathrm{nom}} \cdot \min\!\Bigl( \frac{J_{xx} + \dJ}{J_{xx}},\; 5.0 \Bigr). \label{eq:gain}
\end{equation}

% =====================================================================
% Optional subsection for Section IV (place before \subsection{Ground-truth-free
% online geometry}) -- pre-empts the reviewer/advisor question directly.
% =====================================================================
\subsection{Why the inertia and CoM are not augmented into the state}
\label{sec:whynotaugment}
Augmenting $\dJ$ and $\cvec$ into the MHE decision vector is the obvious alternative, and we
argue against it on three grounds. \emph{(i) Asymmetric observability.} The mass is excited
directly by the thrust--acceleration relation \eqref{eq:trans} and is observable in near-hover
flight; $\dJ$ acts only through $\dot{\boldsymbol{\omega}}$ in \eqref{eq:rot}, which a delivery
task in near-hover excites weakly. Making it free lets the optimizer explain noise with a
near-unobservable quantity. \emph{(ii) They compete for the same residual.} Deriving the
geometry from the \emph{instantaneous} mass estimate closes a feedback loop
$m_{\mathrm{est}} \!\rightarrow\! \dJ,\cvec \!\rightarrow\! \text{attitude model} \!\rightarrow\! m_{\mathrm{est}}$;
we observed this bootstrap fail structurally on light payloads (\S\ref{sec:method}-E), and the
fix was to sever it, not to add more free parameters. \emph{(iii) Economy.} Equations
\eqref{eq:dj}--\eqref{eq:cxy} obtain all three numbers from one scalar $\mpay$, and that scalar
may be a weak operational prior: with the geometry prior set \SI{33}{\percent} wrong, the mass
estimate still converges to within \SI{0.3}{\percent} (\S\ref{sec:results}-C). The online
variant \eqref{eq:conline} then refines $\cvec$ through a channel that reads no model state at
all:
\begin{equation}
c_x = -\,\tau_{\mathrm{pitch}}^{\mathrm{phys}} / \Tphys, \qquad
c_y = +\,\tau_{\mathrm{roll}}^{\mathrm{phys}} / \Tphys, \label{eq:conline}
\end{equation}
gated to steady flight ($|\boldsymbol{\omega}| < \SI{0.15}{\radian\per\second}$,
$|\mathbf{v}_{xy}| < \SI{0.20}{\meter\per\second}$) and low-pass filtered
($\tau = \SI{2}{s}$). Both $\Tphys$ and $\tauphys$ come from rotor speeds
($F_i = k_f \omega_i^2$), so \eqref{eq:conline} is an algebraic inversion, not an optimization,
and is structurally decoupled from $m$.

% =====================================================================
% Two sentences for Section VII (Limitations) -- declare the one modelling
% approximation that is genuinely non-negligible, before a reviewer finds it.
% =====================================================================
\textbf{[Limitations insert]} Two approximations in \eqref{eq:dj} are worth stating. The
single scalar $\dJ$ replaces the pair $(\dJ_{xx}, \dJ_{yy})$ by their mean, an anisotropy
error of $\tfrac12 \mu |r_x^2 - r_y^2|$, below \SI{1}{\percent} of $J_{xx}+\dJ$ at our offsets.
More substantially, the yaw increment $\dJ_{zz} = \mu(r_x^2 + r_y^2)$ is neglected: at
\SI{0.3}{\kg} and $r_{xy} = \SI{0.109}{\meter}$ it is \SI{0.0031}{\kg\meter\squared}, some
\SI{15}{\percent} of $J_{zz}$. The yaw channel neither carries the payload compensation nor
participates in \eqref{eq:conline}, so this does not affect the reported results; removing the
approximation requires only replacing $J_{zz}$ by $J_{zz} + \mu(r_x^2+r_y^2)$ in
\eqref{eq:rot} and introduces no new estimated quantity.

%%   (b) Standalone preview:  latexmk -pdf preview_sec3.tex
%%
%% REQUIRES in the preamble (all already present in main.tex):
%%   amsmath, booktabs, siunitx, multirow, array
%%
%% NOTE ON NOTATION: main.tex currently writes the inertia increment as
%% "dJ" (\S III, \S IV-E, \S VI-C). This file uses \dJ -> \Delta J, which
%% is unambiguous against the differential d. If you adopt it, sweep the
%% remaining occurrences: grep -n "dJ" main.tex
%% =====================================================================

% ---- Notation macros (safe to move to the preamble) -----------------
\providecommand{\dJ}{\Delta J}
\providecommand{\cvec}{\mathbf{c}}
\providecommand{\rp}{\mathbf{r}_p}
\providecommand{\mpay}{\hat{m}_p}          % payload-mass scalar (prior)
\providecommand{\mb}{m_b}                % bare-airframe mass
\providecommand{\mt}{m_t}                % composite mass
\providecommand{\Tphys}{T_{\mathrm{phys}}}
\providecommand{\tauphys}{\boldsymbol{\tau}_{\mathrm{phys}}}
\providecommand{\tcom}{\boldsymbol{\tau}_{\mathrm{com}}}

% =====================================================================
% TABLE I -- parameter identity ledger
% Place early in Section III. Uses table* (two-column span); if space is
% tight, switch to table and drop the "Value / source" column.
% =====================================================================
\begin{table*}[t]
\caption{Model parameters: which are estimated online, which are algebraically
derived, and which are fixed. \emph{Only the mass is a decision variable of the
estimator}; the inertia increment and CoM offset enter the MHE as \emph{known}
runtime parameters, on the same footing as the measured thrust and torque.}
\label{tab:params}
\centering
\footnotesize
\begin{tabular}{@{}l p{0.17\textwidth} c p{0.25\textwidth} p{0.28\textwidth}@{}}
\toprule
Symbol & Meaning & Dim. & Identity & Enters the model through \\
\midrule
$m$ & composite mass (airframe + payload) & 1 & \textbf{MHE decision variable} ($\dot m = 0$, state 14) & translational dynamics \eqref{eq:trans}; cost hover baseline \eqref{eq:hover} \\
$\dJ$ & roll/pitch inertia increment & 1 & algebraically derived, \eqref{eq:dj} & rotational dynamics \eqref{eq:rot}; inner-loop gain scaling \eqref{eq:gain} \\
$\cvec = [c_x, c_y]$ & composite-CoM horizontal offset & 2 & derived \eqref{eq:cxy}, or algebraically inverted online \eqref{eq:conline} & thrust eccentricity torque \eqref{eq:tcom}; cost torque baseline \eqref{eq:hover} \\
\midrule
$\mpay$ & payload-mass scalar (\emph{the single unknown behind $\dJ$ and $\cvec$}) & 1 & weak operational prior (nominal parcel mass) & \eqref{eq:dj}, \eqref{eq:cxy} \\
$\rp = [r_x, r_y, r_z]$ & payload offset in body frame & 3 & known (gripper contact geometry) & \eqref{eq:dj}, \eqref{eq:cxy} \\
\midrule
$\mb$ & bare-airframe mass & 1 & fixed calibration & \SI{2.0643}{\kg} \\
$J_b$ & bare-airframe inertia & 3 & fixed calibration & $\mathrm{diag}(0.0142, 0.0142, 0.0210)$ \si{\kg\meter\squared} \\
$g$ & gravity & 1 & fixed constant & \SI{9.81}{\meter\per\second\squared} \\
\bottomrule
\end{tabular}
\end{table*}

% =====================================================================
% Replacement for the "Dynamics and parameterization" paragraph
% =====================================================================
\textbf{Dynamics and parameterization.} With state
$x = [\mathbf{p}; \mathbf{v}; \mathbf{q}; \boldsymbol{\omega}] \in \mathbb{R}^{13}$ and input
$u = [T; \boldsymbol{\tau}] \in \mathbb{R}^{4}$ expressed in an ENU world frame, and with
$\mathbf{e}_3 = [0, 0, 1]^{\!\top}$ the vertical basis vector, the model is
\begin{align}
\dot{\mathbf{p}} &= \mathbf{v}, \label{eq:kin}\\
\dot{\mathbf{v}} &= \frac{1}{m}\, R(\mathbf{q})\,\mathbf{e}_3 T - g\,\mathbf{e}_3, \label{eq:trans}\\
\dot{\boldsymbol{\omega}} &= J(\dJ)^{-1}\Bigl( \boldsymbol{\tau} + \tcom(\cvec) - \boldsymbol{\omega} \times J(\dJ)\boldsymbol{\omega} \Bigr), \label{eq:rot}
\end{align}
with the attitude kinematics, for $\mathbf{q} = [q_w, q_x, q_y, q_z]^{\!\top}$,
\begin{equation}
\dot{\mathbf{q}} = \frac{1}{2}
\begin{bmatrix}
-q_x & -q_y & -q_z \\
 q_w & -q_z &  q_y \\
 q_z &  q_w & -q_x \\
-q_y &  q_x &  q_w
\end{bmatrix} \boldsymbol{\omega},
\label{eq:quat}
\end{equation}
and $J(\dJ) = \mathrm{diag}\bigl(J_{xx} + \dJ,\; J_{yy} + \dJ,\; J_{zz}\bigr)$. Thrust acts at
the geometric center along the body $z$ axis, so about a composite CoM at
$\mathbf{r}_c = [c_x, c_y, c_z]$ it produces
\begin{equation}
\tcom(\cvec) = (-\mathbf{r}_c) \times \mathbf{e}_3 T = \bigl[-c_y T,\; +c_x T,\; 0\bigr]^{\!\top}. \label{eq:tcom}
\end{equation}
Note that $c_z$ \emph{cannot} appear: a vertical CoM offset is collinear with the body-$z$
thrust and their cross product vanishes identically, and gravity acts at the CoM itself. The
CoM offset is therefore a two-parameter quantity by construction, not by truncation.

Each parameter has exactly one point of entry: $m$ appears only as the thrust-to-acceleration
gain $1/m$ in \eqref{eq:trans}, $\dJ$ only in $J(\dJ)$ in \eqref{eq:rot}, and $\cvec$ only in
\eqref{eq:tcom}. This is what makes $m$ observable from a model-independent thrust
reconstruction: given $\Tphys$, the measured $\dot{\mathbf{v}}$ pins $m$ uniquely.

\textbf{Structured parameterization of the payload geometry.} $\dJ$ and $\cvec$ are
\emph{not} additional estimated degrees of freedom. Modelling the payload as a point mass
$\mpay$ rigidly attached at a known body-frame offset $\rp$, and writing
$\mt = \mb + \mpay$ and the reduced mass $\mu = \mb\,\mpay / \mt$, the parallel-axis theorem and
the CoM definition give them in closed form:
\begin{align}
\dJ &= \tfrac12\bigl(\dJ_{xx} + \dJ_{yy}\bigr) = \mu\Bigl( r_z^2 + \tfrac12\bigl(r_x^2 + r_y^2\bigr) \Bigr), \label{eq:dj}\\
\cvec &= \frac{\mpay}{\mt}\,\bigl[r_x,\; r_y\bigr]. \label{eq:cxy}
\end{align}
Three numbers are thus generated from a \emph{single} unknown scalar $\mpay$ plus known
geometry: the rigid-body constraint replaces two degrees of freedom that a naive
augmentation would have to identify. For a \SI{0.3}{\kg} payload at
$\rp = [0, 0.109, -0.47]\,$\si{\meter} this yields $\mu = \SI{0.262}{\kg}$,
$\dJ = \SI{0.0594}{\kg\meter\squared}$ --- an inertia ratio $(J_{xx}+\dJ)/J_{xx} = 5.2$ ---
and $c_y = \SI{13.8}{\milli\meter}$, i.e.\ a constant eccentric hover torque
$|c_y| \mt g \approx \SI{0.32}{\newton\meter}$, about \SI{64}{\percent} of the available
torque authority. Compensating mass alone is therefore insufficient.

\textbf{MHE.} The estimator augments the mass into the state,
$x_{\mathrm{aug}} = [x;\, m] \in \mathbb{R}^{14}$ with $\dot m = 0$ (constant over the window)
and process noise on the 13 physical states only; the payload geometry is passed in as a
\emph{known} runtime parameter vector alongside the measured input,
\begin{equation}
p_{\mathrm{MHE}} = \bigl[\, T,\; \boldsymbol{\tau},\; \dJ,\; c_x,\; c_y \,\bigr] \in \mathbb{R}^{7}. \label{eq:pmhe}
\end{equation}
Window $N = 20$, $dt = \SI{0.1}{s}$, terminal measurement excluded from the cost, and $m$
bounded to $[\SI{1.961}{\kg}, \SI{5.0}{\kg}]$ (a physical floor at $0.95\mb$: the vehicle
cannot be lighter than the bare airframe).

\textbf{NMPC parameter consumption.} The controller carries
$p_{\mathrm{NMPC}} = [x_{\mathrm{ref}}(13),\, m,\, \dJ,\, \cvec(2)] \in \mathbb{R}^{17}$.
The hover baseline in the input cost is re-centred by the live parameters,
\begin{equation}
T_h = m\,g, \qquad
u_{\mathrm{hover}} = \bigl[\, T_h,\; c_y T_h,\; -c_x T_h,\; 0 \,\bigr]^{\!\top}, \label{eq:hover}
\end{equation}
the torque entries being exactly $-\tcom$ of \eqref{eq:tcom}: hovering with an eccentric load
\emph{requires} a constant trim torque, and a cost that does not admit this pulls the torque
toward zero, producing the same class of steady-state bias as a stale hover-thrust baseline.
Finally, model consistency cannot restore inner-loop bandwidth, so the PX4 rate gains are
scaled at attach by the inertia ratio
\begin{equation}
K \leftarrow K_{\mathrm{nom}} \cdot \min\!\Bigl( \frac{J_{xx} + \dJ}{J_{xx}},\; 5.0 \Bigr). \label{eq:gain}
\end{equation}

% =====================================================================
% Optional subsection for Section IV (place before \subsection{Ground-truth-free
% online geometry}) -- pre-empts the reviewer/advisor question directly.
% =====================================================================
\subsection{Why the inertia and CoM are not augmented into the state}
\label{sec:whynotaugment}
Augmenting $\dJ$ and $\cvec$ into the MHE decision vector is the obvious alternative, and we
argue against it on three grounds. \emph{(i) Asymmetric observability.} The mass is excited
directly by the thrust--acceleration relation \eqref{eq:trans} and is observable in near-hover
flight; $\dJ$ acts only through $\dot{\boldsymbol{\omega}}$ in \eqref{eq:rot}, which a delivery
task in near-hover excites weakly. Making it free lets the optimizer explain noise with a
near-unobservable quantity. \emph{(ii) They compete for the same residual.} Deriving the
geometry from the \emph{instantaneous} mass estimate closes a feedback loop
$m_{\mathrm{est}} \!\rightarrow\! \dJ,\cvec \!\rightarrow\! \text{attitude model} \!\rightarrow\! m_{\mathrm{est}}$;
we observed this bootstrap fail structurally on light payloads (\S\ref{sec:method}-E), and the
fix was to sever it, not to add more free parameters. \emph{(iii) Economy.} Equations
\eqref{eq:dj}--\eqref{eq:cxy} obtain all three numbers from one scalar $\mpay$, and that scalar
may be a weak operational prior: with the geometry prior set \SI{33}{\percent} wrong, the mass
estimate still converges to within \SI{0.3}{\percent} (\S\ref{sec:results}-C). The online
variant \eqref{eq:conline} then refines $\cvec$ through a channel that reads no model state at
all:
\begin{equation}
c_x = -\,\tau_{\mathrm{pitch}}^{\mathrm{phys}} / \Tphys, \qquad
c_y = +\,\tau_{\mathrm{roll}}^{\mathrm{phys}} / \Tphys, \label{eq:conline}
\end{equation}
gated to steady flight ($|\boldsymbol{\omega}| < \SI{0.15}{\radian\per\second}$,
$|\mathbf{v}_{xy}| < \SI{0.20}{\meter\per\second}$) and low-pass filtered
($\tau = \SI{2}{s}$). Both $\Tphys$ and $\tauphys$ come from rotor speeds
($F_i = k_f \omega_i^2$), so \eqref{eq:conline} is an algebraic inversion, not an optimization,
and is structurally decoupled from $m$.

% =====================================================================
% Two sentences for Section VII (Limitations) -- declare the one modelling
% approximation that is genuinely non-negligible, before a reviewer finds it.
% =====================================================================
\textbf{[Limitations insert]} Two approximations in \eqref{eq:dj} are worth stating. The
single scalar $\dJ$ replaces the pair $(\dJ_{xx}, \dJ_{yy})$ by their mean, an anisotropy
error of $\tfrac12 \mu |r_x^2 - r_y^2|$, below \SI{1}{\percent} of $J_{xx}+\dJ$ at our offsets.
More substantially, the yaw increment $\dJ_{zz} = \mu(r_x^2 + r_y^2)$ is neglected: at
\SI{0.3}{\kg} and $r_{xy} = \SI{0.109}{\meter}$ it is \SI{0.0031}{\kg\meter\squared}, some
\SI{15}{\percent} of $J_{zz}$. The yaw channel neither carries the payload compensation nor
participates in \eqref{eq:conline}, so this does not affect the reported results; removing the
approximation requires only replacing $J_{zz}$ by $J_{zz} + \mu(r_x^2+r_y^2)$ in
\eqref{eq:rot} and introduces no new estimated quantity.

%%   (b) Standalone preview:  latexmk -pdf preview_sec3.tex
%%
%% REQUIRES in the preamble (all already present in main.tex):
%%   amsmath, booktabs, siunitx, multirow, array
%%
%% NOTE ON NOTATION: main.tex currently writes the inertia increment as
%% "dJ" (\S III, \S IV-E, \S VI-C). This file uses \dJ -> \Delta J, which
%% is unambiguous against the differential d. If you adopt it, sweep the
%% remaining occurrences: grep -n "dJ" main.tex
%% =====================================================================

% ---- Notation macros (safe to move to the preamble) -----------------
\providecommand{\dJ}{\Delta J}
\providecommand{\cvec}{\mathbf{c}}
\providecommand{\rp}{\mathbf{r}_p}
\providecommand{\mpay}{\hat{m}_p}          % payload-mass scalar (prior)
\providecommand{\mb}{m_b}                % bare-airframe mass
\providecommand{\mt}{m_t}                % composite mass
\providecommand{\Tphys}{T_{\mathrm{phys}}}
\providecommand{\tauphys}{\boldsymbol{\tau}_{\mathrm{phys}}}
\providecommand{\tcom}{\boldsymbol{\tau}_{\mathrm{com}}}

% =====================================================================
% TABLE I -- parameter identity ledger
% Place early in Section III. Uses table* (two-column span); if space is
% tight, switch to table and drop the "Value / source" column.
% =====================================================================
\begin{table*}[t]
\caption{Model parameters: which are estimated online, which are algebraically
derived, and which are fixed. \emph{Only the mass is a decision variable of the
estimator}; the inertia increment and CoM offset enter the MHE as \emph{known}
runtime parameters, on the same footing as the measured thrust and torque.}
\label{tab:params}
\centering
\footnotesize
\begin{tabular}{@{}l p{0.17\textwidth} c p{0.25\textwidth} p{0.28\textwidth}@{}}
\toprule
Symbol & Meaning & Dim. & Identity & Enters the model through \\
\midrule
$m$ & composite mass (airframe + payload) & 1 & \textbf{MHE decision variable} ($\dot m = 0$, state 14) & translational dynamics \eqref{eq:trans}; cost hover baseline \eqref{eq:hover} \\
$\dJ$ & roll/pitch inertia increment & 1 & algebraically derived, \eqref{eq:dj} & rotational dynamics \eqref{eq:rot}; inner-loop gain scaling \eqref{eq:gain} \\
$\cvec = [c_x, c_y]$ & composite-CoM horizontal offset & 2 & derived \eqref{eq:cxy}, or algebraically inverted online \eqref{eq:conline} & thrust eccentricity torque \eqref{eq:tcom}; cost torque baseline \eqref{eq:hover} \\
\midrule
$\mpay$ & payload-mass scalar (\emph{the single unknown behind $\dJ$ and $\cvec$}) & 1 & weak operational prior (nominal parcel mass) & \eqref{eq:dj}, \eqref{eq:cxy} \\
$\rp = [r_x, r_y, r_z]$ & payload offset in body frame & 3 & known (gripper contact geometry) & \eqref{eq:dj}, \eqref{eq:cxy} \\
\midrule
$\mb$ & bare-airframe mass & 1 & fixed calibration & \SI{2.0643}{\kg} \\
$J_b$ & bare-airframe inertia & 3 & fixed calibration & $\mathrm{diag}(0.0142, 0.0142, 0.0210)$ \si{\kg\meter\squared} \\
$g$ & gravity & 1 & fixed constant & \SI{9.81}{\meter\per\second\squared} \\
\bottomrule
\end{tabular}
\end{table*}

% =====================================================================
% Replacement for the "Dynamics and parameterization" paragraph
% =====================================================================
\textbf{Dynamics and parameterization.} With state
$x = [\mathbf{p}; \mathbf{v}; \mathbf{q}; \boldsymbol{\omega}] \in \mathbb{R}^{13}$ and input
$u = [T; \boldsymbol{\tau}] \in \mathbb{R}^{4}$ expressed in an ENU world frame, and with
$\mathbf{e}_3 = [0, 0, 1]^{\!\top}$ the vertical basis vector, the model is
\begin{align}
\dot{\mathbf{p}} &= \mathbf{v}, \label{eq:kin}\\
\dot{\mathbf{v}} &= \frac{1}{m}\, R(\mathbf{q})\,\mathbf{e}_3 T - g\,\mathbf{e}_3, \label{eq:trans}\\
\dot{\boldsymbol{\omega}} &= J(\dJ)^{-1}\Bigl( \boldsymbol{\tau} + \tcom(\cvec) - \boldsymbol{\omega} \times J(\dJ)\boldsymbol{\omega} \Bigr), \label{eq:rot}
\end{align}
with the attitude kinematics, for $\mathbf{q} = [q_w, q_x, q_y, q_z]^{\!\top}$,
\begin{equation}
\dot{\mathbf{q}} = \frac{1}{2}
\begin{bmatrix}
-q_x & -q_y & -q_z \\
 q_w & -q_z &  q_y \\
 q_z &  q_w & -q_x \\
-q_y &  q_x &  q_w
\end{bmatrix} \boldsymbol{\omega},
\label{eq:quat}
\end{equation}
and $J(\dJ) = \mathrm{diag}\bigl(J_{xx} + \dJ,\; J_{yy} + \dJ,\; J_{zz}\bigr)$. Thrust acts at
the geometric center along the body $z$ axis, so about a composite CoM at
$\mathbf{r}_c = [c_x, c_y, c_z]$ it produces
\begin{equation}
\tcom(\cvec) = (-\mathbf{r}_c) \times \mathbf{e}_3 T = \bigl[-c_y T,\; +c_x T,\; 0\bigr]^{\!\top}. \label{eq:tcom}
\end{equation}
Note that $c_z$ \emph{cannot} appear: a vertical CoM offset is collinear with the body-$z$
thrust and their cross product vanishes identically, and gravity acts at the CoM itself. The
CoM offset is therefore a two-parameter quantity by construction, not by truncation.

Each parameter has exactly one point of entry: $m$ appears only as the thrust-to-acceleration
gain $1/m$ in \eqref{eq:trans}, $\dJ$ only in $J(\dJ)$ in \eqref{eq:rot}, and $\cvec$ only in
\eqref{eq:tcom}. This is what makes $m$ observable from a model-independent thrust
reconstruction: given $\Tphys$, the measured $\dot{\mathbf{v}}$ pins $m$ uniquely.

\textbf{Structured parameterization of the payload geometry.} $\dJ$ and $\cvec$ are
\emph{not} additional estimated degrees of freedom. Modelling the payload as a point mass
$\mpay$ rigidly attached at a known body-frame offset $\rp$, and writing
$\mt = \mb + \mpay$ and the reduced mass $\mu = \mb\,\mpay / \mt$, the parallel-axis theorem and
the CoM definition give them in closed form:
\begin{align}
\dJ &= \tfrac12\bigl(\dJ_{xx} + \dJ_{yy}\bigr) = \mu\Bigl( r_z^2 + \tfrac12\bigl(r_x^2 + r_y^2\bigr) \Bigr), \label{eq:dj}\\
\cvec &= \frac{\mpay}{\mt}\,\bigl[r_x,\; r_y\bigr]. \label{eq:cxy}
\end{align}
Three numbers are thus generated from a \emph{single} unknown scalar $\mpay$ plus known
geometry: the rigid-body constraint replaces two degrees of freedom that a naive
augmentation would have to identify. For a \SI{0.3}{\kg} payload at
$\rp = [0, 0.109, -0.47]\,$\si{\meter} this yields $\mu = \SI{0.262}{\kg}$,
$\dJ = \SI{0.0594}{\kg\meter\squared}$ --- an inertia ratio $(J_{xx}+\dJ)/J_{xx} = 5.2$ ---
and $c_y = \SI{13.8}{\milli\meter}$, i.e.\ a constant eccentric hover torque
$|c_y| \mt g \approx \SI{0.32}{\newton\meter}$, about \SI{64}{\percent} of the available
torque authority. Compensating mass alone is therefore insufficient.

\textbf{MHE.} The estimator augments the mass into the state,
$x_{\mathrm{aug}} = [x;\, m] \in \mathbb{R}^{14}$ with $\dot m = 0$ (constant over the window)
and process noise on the 13 physical states only; the payload geometry is passed in as a
\emph{known} runtime parameter vector alongside the measured input,
\begin{equation}
p_{\mathrm{MHE}} = \bigl[\, T,\; \boldsymbol{\tau},\; \dJ,\; c_x,\; c_y \,\bigr] \in \mathbb{R}^{7}. \label{eq:pmhe}
\end{equation}
Window $N = 20$, $dt = \SI{0.1}{s}$, terminal measurement excluded from the cost, and $m$
bounded to $[\SI{1.961}{\kg}, \SI{5.0}{\kg}]$ (a physical floor at $0.95\mb$: the vehicle
cannot be lighter than the bare airframe).

\textbf{NMPC parameter consumption.} The controller carries
$p_{\mathrm{NMPC}} = [x_{\mathrm{ref}}(13),\, m,\, \dJ,\, \cvec(2)] \in \mathbb{R}^{17}$.
The hover baseline in the input cost is re-centred by the live parameters,
\begin{equation}
T_h = m\,g, \qquad
u_{\mathrm{hover}} = \bigl[\, T_h,\; c_y T_h,\; -c_x T_h,\; 0 \,\bigr]^{\!\top}, \label{eq:hover}
\end{equation}
the torque entries being exactly $-\tcom$ of \eqref{eq:tcom}: hovering with an eccentric load
\emph{requires} a constant trim torque, and a cost that does not admit this pulls the torque
toward zero, producing the same class of steady-state bias as a stale hover-thrust baseline.
Finally, model consistency cannot restore inner-loop bandwidth, so the PX4 rate gains are
scaled at attach by the inertia ratio
\begin{equation}
K \leftarrow K_{\mathrm{nom}} \cdot \min\!\Bigl( \frac{J_{xx} + \dJ}{J_{xx}},\; 5.0 \Bigr). \label{eq:gain}
\end{equation}

% =====================================================================
% Optional subsection for Section IV (place before \subsection{Ground-truth-free
% online geometry}) -- pre-empts the reviewer/advisor question directly.
% =====================================================================
\subsection{Why the inertia and CoM are not augmented into the state}
\label{sec:whynotaugment}
Augmenting $\dJ$ and $\cvec$ into the MHE decision vector is the obvious alternative, and we
argue against it on three grounds. \emph{(i) Asymmetric observability.} The mass is excited
directly by the thrust--acceleration relation \eqref{eq:trans} and is observable in near-hover
flight; $\dJ$ acts only through $\dot{\boldsymbol{\omega}}$ in \eqref{eq:rot}, which a delivery
task in near-hover excites weakly. Making it free lets the optimizer explain noise with a
near-unobservable quantity. \emph{(ii) They compete for the same residual.} Deriving the
geometry from the \emph{instantaneous} mass estimate closes a feedback loop
$m_{\mathrm{est}} \!\rightarrow\! \dJ,\cvec \!\rightarrow\! \text{attitude model} \!\rightarrow\! m_{\mathrm{est}}$;
we observed this bootstrap fail structurally on light payloads (\S\ref{sec:method}-E), and the
fix was to sever it, not to add more free parameters. \emph{(iii) Economy.} Equations
\eqref{eq:dj}--\eqref{eq:cxy} obtain all three numbers from one scalar $\mpay$, and that scalar
may be a weak operational prior: with the geometry prior set \SI{33}{\percent} wrong, the mass
estimate still converges to within \SI{0.3}{\percent} (\S\ref{sec:results}-C). The online
variant \eqref{eq:conline} then refines $\cvec$ through a channel that reads no model state at
all:
\begin{equation}
c_x = -\,\tau_{\mathrm{pitch}}^{\mathrm{phys}} / \Tphys, \qquad
c_y = +\,\tau_{\mathrm{roll}}^{\mathrm{phys}} / \Tphys, \label{eq:conline}
\end{equation}
gated to steady flight ($|\boldsymbol{\omega}| < \SI{0.15}{\radian\per\second}$,
$|\mathbf{v}_{xy}| < \SI{0.20}{\meter\per\second}$) and low-pass filtered
($\tau = \SI{2}{s}$). Both $\Tphys$ and $\tauphys$ come from rotor speeds
($F_i = k_f \omega_i^2$), so \eqref{eq:conline} is an algebraic inversion, not an optimization,
and is structurally decoupled from $m$.

% =====================================================================
% Two sentences for Section VII (Limitations) -- declare the one modelling
% approximation that is genuinely non-negligible, before a reviewer finds it.
% =====================================================================
\textbf{[Limitations insert]} Two approximations in \eqref{eq:dj} are worth stating. The
single scalar $\dJ$ replaces the pair $(\dJ_{xx}, \dJ_{yy})$ by their mean, an anisotropy
error of $\tfrac12 \mu |r_x^2 - r_y^2|$, below \SI{1}{\percent} of $J_{xx}+\dJ$ at our offsets.
More substantially, the yaw increment $\dJ_{zz} = \mu(r_x^2 + r_y^2)$ is neglected: at
\SI{0.3}{\kg} and $r_{xy} = \SI{0.109}{\meter}$ it is \SI{0.0031}{\kg\meter\squared}, some
\SI{15}{\percent} of $J_{zz}$. The yaw channel neither carries the payload compensation nor
participates in \eqref{eq:conline}, so this does not affect the reported results; removing the
approximation requires only replacing $J_{zz}$ by $J_{zz} + \mu(r_x^2+r_y^2)$ in
\eqref{eq:rot} and introduces no new estimated quantity.

%%   (b) Standalone preview:  latexmk -pdf preview_sec3.tex
%%
%% REQUIRES in the preamble (all already present in main.tex):
%%   amsmath, booktabs, siunitx, multirow, array
%%
%% NOTE ON NOTATION: main.tex currently writes the inertia increment as
%% "dJ" (\S III, \S IV-E, \S VI-C). This file uses \dJ -> \Delta J, which
%% is unambiguous against the differential d. If you adopt it, sweep the
%% remaining occurrences: grep -n "dJ" main.tex
%% =====================================================================

% ---- Notation macros (safe to move to the preamble) -----------------
\providecommand{\dJ}{\Delta J}
\providecommand{\cvec}{\mathbf{c}}
\providecommand{\rp}{\mathbf{r}_p}
\providecommand{\mpay}{\hat{m}_p}          % payload-mass scalar (prior)
\providecommand{\mb}{m_b}                % bare-airframe mass
\providecommand{\mt}{m_t}                % composite mass
\providecommand{\Tphys}{T_{\mathrm{phys}}}
\providecommand{\tauphys}{\boldsymbol{\tau}_{\mathrm{phys}}}
\providecommand{\tcom}{\boldsymbol{\tau}_{\mathrm{com}}}

% =====================================================================
% TABLE I -- parameter identity ledger
% Place early in Section III. Uses table* (two-column span); if space is
% tight, switch to table and drop the "Value / source" column.
% =====================================================================
\begin{table*}[t]
\caption{Model parameters: which are estimated online, which are algebraically
derived, and which are fixed. \emph{Only the mass is a decision variable of the
estimator}; the inertia increment and CoM offset enter the MHE as \emph{known}
runtime parameters, on the same footing as the measured thrust and torque.}
\label{tab:params}
\centering
\footnotesize
\begin{tabular}{@{}l p{0.17\textwidth} c p{0.25\textwidth} p{0.28\textwidth}@{}}
\toprule
Symbol & Meaning & Dim. & Identity & Enters the model through \\
\midrule
$m$ & composite mass (airframe + payload) & 1 & \textbf{MHE decision variable} ($\dot m = 0$, state 14) & translational dynamics \eqref{eq:trans}; cost hover baseline \eqref{eq:hover} \\
$\dJ$ & roll/pitch inertia increment & 1 & algebraically derived, \eqref{eq:dj} & rotational dynamics \eqref{eq:rot}; inner-loop gain scaling \eqref{eq:gain} \\
$\cvec = [c_x, c_y]$ & composite-CoM horizontal offset & 2 & derived \eqref{eq:cxy}, or algebraically inverted online \eqref{eq:conline} & thrust eccentricity torque \eqref{eq:tcom}; cost torque baseline \eqref{eq:hover} \\
\midrule
$\mpay$ & payload-mass scalar (\emph{the single unknown behind $\dJ$ and $\cvec$}) & 1 & weak operational prior (nominal parcel mass) & \eqref{eq:dj}, \eqref{eq:cxy} \\
$\rp = [r_x, r_y, r_z]$ & payload offset in body frame & 3 & known (gripper contact geometry) & \eqref{eq:dj}, \eqref{eq:cxy} \\
\midrule
$\mb$ & bare-airframe mass & 1 & fixed calibration & \SI{2.0643}{\kg} \\
$J_b$ & bare-airframe inertia & 3 & fixed calibration & $\mathrm{diag}(0.0142, 0.0142, 0.0210)$ \si{\kg\meter\squared} \\
$g$ & gravity & 1 & fixed constant & \SI{9.81}{\meter\per\second\squared} \\
\bottomrule
\end{tabular}
\end{table*}

% =====================================================================
% Replacement for the "Dynamics and parameterization" paragraph
% =====================================================================
\textbf{Dynamics and parameterization.} With state
$x = [\mathbf{p}; \mathbf{v}; \mathbf{q}; \boldsymbol{\omega}] \in \mathbb{R}^{13}$ and input
$u = [T; \boldsymbol{\tau}] \in \mathbb{R}^{4}$ expressed in an ENU world frame, and with
$\mathbf{e}_3 = [0, 0, 1]^{\!\top}$ the vertical basis vector, the model is
\begin{align}
\dot{\mathbf{p}} &= \mathbf{v}, \label{eq:kin}\\
\dot{\mathbf{v}} &= \frac{1}{m}\, R(\mathbf{q})\,\mathbf{e}_3 T - g\,\mathbf{e}_3, \label{eq:trans}\\
\dot{\boldsymbol{\omega}} &= J(\dJ)^{-1}\Bigl( \boldsymbol{\tau} + \tcom(\cvec) - \boldsymbol{\omega} \times J(\dJ)\boldsymbol{\omega} \Bigr), \label{eq:rot}
\end{align}
with the attitude kinematics, for $\mathbf{q} = [q_w, q_x, q_y, q_z]^{\!\top}$,
\begin{equation}
\dot{\mathbf{q}} = \frac{1}{2}
\begin{bmatrix}
-q_x & -q_y & -q_z \\
 q_w & -q_z &  q_y \\
 q_z &  q_w & -q_x \\
-q_y &  q_x &  q_w
\end{bmatrix} \boldsymbol{\omega},
\label{eq:quat}
\end{equation}
and $J(\dJ) = \mathrm{diag}\bigl(J_{xx} + \dJ,\; J_{yy} + \dJ,\; J_{zz}\bigr)$. Thrust acts at
the geometric center along the body $z$ axis, so about a composite CoM at
$\mathbf{r}_c = [c_x, c_y, c_z]$ it produces
\begin{equation}
\tcom(\cvec) = (-\mathbf{r}_c) \times \mathbf{e}_3 T = \bigl[-c_y T,\; +c_x T,\; 0\bigr]^{\!\top}. \label{eq:tcom}
\end{equation}
Note that $c_z$ \emph{cannot} appear: a vertical CoM offset is collinear with the body-$z$
thrust and their cross product vanishes identically, and gravity acts at the CoM itself. The
CoM offset is therefore a two-parameter quantity by construction, not by truncation.

Each parameter has exactly one point of entry: $m$ appears only as the thrust-to-acceleration
gain $1/m$ in \eqref{eq:trans}, $\dJ$ only in $J(\dJ)$ in \eqref{eq:rot}, and $\cvec$ only in
\eqref{eq:tcom}. This is what makes $m$ observable from a model-independent thrust
reconstruction: given $\Tphys$, the measured $\dot{\mathbf{v}}$ pins $m$ uniquely.

\textbf{Structured parameterization of the payload geometry.} $\dJ$ and $\cvec$ are
\emph{not} additional estimated degrees of freedom. Modelling the payload as a point mass
$\mpay$ rigidly attached at a known body-frame offset $\rp$, and writing
$\mt = \mb + \mpay$ and the reduced mass $\mu = \mb\,\mpay / \mt$, the parallel-axis theorem and
the CoM definition give them in closed form:
\begin{align}
\dJ &= \tfrac12\bigl(\dJ_{xx} + \dJ_{yy}\bigr) = \mu\Bigl( r_z^2 + \tfrac12\bigl(r_x^2 + r_y^2\bigr) \Bigr), \label{eq:dj}\\
\cvec &= \frac{\mpay}{\mt}\,\bigl[r_x,\; r_y\bigr]. \label{eq:cxy}
\end{align}
Three numbers are thus generated from a \emph{single} unknown scalar $\mpay$ plus known
geometry: the rigid-body constraint replaces two degrees of freedom that a naive
augmentation would have to identify. For a \SI{0.3}{\kg} payload at
$\rp = [0, 0.109, -0.47]\,$\si{\meter} this yields $\mu = \SI{0.262}{\kg}$,
$\dJ = \SI{0.0594}{\kg\meter\squared}$ --- an inertia ratio $(J_{xx}+\dJ)/J_{xx} = 5.2$ ---
and $c_y = \SI{13.8}{\milli\meter}$, i.e.\ a constant eccentric hover torque
$|c_y| \mt g \approx \SI{0.32}{\newton\meter}$, about \SI{64}{\percent} of the available
torque authority. Compensating mass alone is therefore insufficient.

\textbf{MHE.} The estimator augments the mass into the state,
$x_{\mathrm{aug}} = [x;\, m] \in \mathbb{R}^{14}$ with $\dot m = 0$ (constant over the window)
and process noise on the 13 physical states only; the payload geometry is passed in as a
\emph{known} runtime parameter vector alongside the measured input,
\begin{equation}
p_{\mathrm{MHE}} = \bigl[\, T,\; \boldsymbol{\tau},\; \dJ,\; c_x,\; c_y \,\bigr] \in \mathbb{R}^{7}. \label{eq:pmhe}
\end{equation}
Window $N = 20$, $dt = \SI{0.1}{s}$, terminal measurement excluded from the cost, and $m$
bounded to $[\SI{1.961}{\kg}, \SI{5.0}{\kg}]$ (a physical floor at $0.95\mb$: the vehicle
cannot be lighter than the bare airframe).

\textbf{NMPC parameter consumption.} The controller carries
$p_{\mathrm{NMPC}} = [x_{\mathrm{ref}}(13),\, m,\, \dJ,\, \cvec(2)] \in \mathbb{R}^{17}$.
The hover baseline in the input cost is re-centred by the live parameters,
\begin{equation}
T_h = m\,g, \qquad
u_{\mathrm{hover}} = \bigl[\, T_h,\; c_y T_h,\; -c_x T_h,\; 0 \,\bigr]^{\!\top}, \label{eq:hover}
\end{equation}
the torque entries being exactly $-\tcom$ of \eqref{eq:tcom}: hovering with an eccentric load
\emph{requires} a constant trim torque, and a cost that does not admit this pulls the torque
toward zero, producing the same class of steady-state bias as a stale hover-thrust baseline.
Finally, model consistency cannot restore inner-loop bandwidth, so the PX4 rate gains are
scaled at attach by the inertia ratio
\begin{equation}
K \leftarrow K_{\mathrm{nom}} \cdot \min\!\Bigl( \frac{J_{xx} + \dJ}{J_{xx}},\; 5.0 \Bigr). \label{eq:gain}
\end{equation}

% =====================================================================
% Optional subsection for Section IV (place before \subsection{Ground-truth-free
% online geometry}) -- pre-empts the reviewer/advisor question directly.
% =====================================================================
\subsection{Why the inertia and CoM are not augmented into the state}
\label{sec:whynotaugment}
Augmenting $\dJ$ and $\cvec$ into the MHE decision vector is the obvious alternative, and we
argue against it on three grounds. \emph{(i) Asymmetric observability.} The mass is excited
directly by the thrust--acceleration relation \eqref{eq:trans} and is observable in near-hover
flight; $\dJ$ acts only through $\dot{\boldsymbol{\omega}}$ in \eqref{eq:rot}, which a delivery
task in near-hover excites weakly. Making it free lets the optimizer explain noise with a
near-unobservable quantity. \emph{(ii) They compete for the same residual.} Deriving the
geometry from the \emph{instantaneous} mass estimate closes a feedback loop
$m_{\mathrm{est}} \!\rightarrow\! \dJ,\cvec \!\rightarrow\! \text{attitude model} \!\rightarrow\! m_{\mathrm{est}}$;
we observed this bootstrap fail structurally on light payloads (\S\ref{sec:method}-E), and the
fix was to sever it, not to add more free parameters. \emph{(iii) Economy.} Equations
\eqref{eq:dj}--\eqref{eq:cxy} obtain all three numbers from one scalar $\mpay$, and that scalar
may be a weak operational prior: with the geometry prior set \SI{33}{\percent} wrong, the mass
estimate still converges to within \SI{0.3}{\percent} (\S\ref{sec:results}-C). The online
variant \eqref{eq:conline} then refines $\cvec$ through a channel that reads no model state at
all:
\begin{equation}
c_x = -\,\tau_{\mathrm{pitch}}^{\mathrm{phys}} / \Tphys, \qquad
c_y = +\,\tau_{\mathrm{roll}}^{\mathrm{phys}} / \Tphys, \label{eq:conline}
\end{equation}
gated to steady flight ($|\boldsymbol{\omega}| < \SI{0.15}{\radian\per\second}$,
$|\mathbf{v}_{xy}| < \SI{0.20}{\meter\per\second}$) and low-pass filtered
($\tau = \SI{2}{s}$). Both $\Tphys$ and $\tauphys$ come from rotor speeds
($F_i = k_f \omega_i^2$), so \eqref{eq:conline} is an algebraic inversion, not an optimization,
and is structurally decoupled from $m$.

% =====================================================================
% Two sentences for Section VII (Limitations) -- declare the one modelling
% approximation that is genuinely non-negligible, before a reviewer finds it.
% =====================================================================
\textbf{[Limitations insert]} Two approximations in \eqref{eq:dj} are worth stating. The
single scalar $\dJ$ replaces the pair $(\dJ_{xx}, \dJ_{yy})$ by their mean, an anisotropy
error of $\tfrac12 \mu |r_x^2 - r_y^2|$, below \SI{1}{\percent} of $J_{xx}+\dJ$ at our offsets.
More substantially, the yaw increment $\dJ_{zz} = \mu(r_x^2 + r_y^2)$ is neglected: at
\SI{0.3}{\kg} and $r_{xy} = \SI{0.109}{\meter}$ it is \SI{0.0031}{\kg\meter\squared}, some
\SI{15}{\percent} of $J_{zz}$. The yaw channel neither carries the payload compensation nor
participates in \eqref{eq:conline}, so this does not affect the reported results; removing the
approximation requires only replacing $J_{zz}$ by $J_{zz} + \mu(r_x^2+r_y^2)$ in
\eqref{eq:rot} and introduces no new estimated quantity.
